% ═══════════════════════════════════════════════════════════════════════════════
% HIPÓTESIS - INSTITUTO ATENEO DE COLIMA
% Estudio Documental: Corrupción en México (2006-2012)
% ═══════════════════════════════════════════════════════════════════════════════

\nota{
  Este capítulo formula la hipótesis de trabajo conforme a los estándares de 
  Hernández Sampieri (2014) para estudios documentales de casos múltiples. 
  La hipótesis orienta el análisis documental sin pretender generalización 
  estadística.
}

\section{Hipótesis de trabajo u orientadora}

En congruencia con los lineamientos de Hernández Sampieri et al. (2014) para estudios cualitativos y mixtos de modelo dominante cualitativo basados en casos múltiples, en esta investigación se formuló una hipótesis de trabajo u orientadora (Hi) ---y no una hipótesis causal estadística sujeta a prueba probabilística--- con el propósito de guiar la recolección documental, la codificación categorial y la integración de metainferencias:

\textbf{Hipótesis de trabajo u orientadora (Hi):} Los episodios emblemáticos de corrupción registrados durante el sexenio de Felipe Calderón (2006--2012) no respondieron a fallas administrativas aisladas, sino a condiciones estructurales de concentración discrecional de poder sin contrapesos, planeación presupuestal deficiente en obra pública y desregulación de servicios subrogados con conflicto de interés, lo cual se asoció con efectos prolongados sobre la integridad institucional, el erario público y los derechos humanos ---observables en los casos Genaro García Luna, Refinería Bicentenario y Guardería ABC---, cuya magnitud verificable difiere con frecuencia de la jerarquía asignada por el encuadre mediático.

De la hipótesis de trabajo se derivaron tres supuestos orientadores específicos correspondientes a cada caso principal:

\begin{enumerate}
  \item \textbf{Supuesto 1 (Seguridad y captura institucional):} La concentración de mando operativo, inteligencia y presupuesto en la Secretaría de Seguridad Pública, sin contrapesos efectivos de fiscalización interna ni control bilateral, generó condiciones propicias para la captura institucional y la operación paralela de redes de cohecho criminal y desvío contractual.
  
  \item \textbf{Supuesto 2 (Megainfraestructura y finanzas públicas):} La adopción de decisiones políticas prematuras en megainfraestructura, carente de estudios técnicos vinculantes y candados presupuestales transexenales, convirtió proyectos estratégicos como la Refinería Bicentenario en pasivos fiscales prolongados sin rentabilidad social ni productiva.
  
  \item \textbf{Supuesto 3 (Subrogación y derechos humanos):} La expansión acelerada de la subrogación de servicios públicos esenciales bajo requisitos flexibilizados y vínculos de conflicto de interés diluyó la responsabilidad administrativa del Estado y elevó los riesgos sobre derechos fundamentales, como ocurrió en la Guardería ABC.
\end{enumerate}

Se estableció explícitamente que la evidencia recopilada a partir de tres casos intencionales únicamente permite valorar la congruencia empírico-documental con la hipótesis de trabajo y generar hallazgos analíticamente transferibles, sin pretender una comprobación definitiva ni una generalización estadística hacia la totalidad de la Administración Pública Federal 2006--2012 \citep{hernandez2014}.

\section{Operacionalización de categorías y variables}

\begin{table}[H]
\centering
\caption{Operacionalización de categorías, variables e indicadores documentales}
\label{tab:operacionalizacion}
\small
\begin{tabular}{@{}p{3.5cm}p{5cm}p{5cm}@{}}
\toprule
\textbf{Componente analítico} & \textbf{Categorías / variables examinadas} & \textbf{Indicadores empíricos y documentales verificados} \\
\midrule
\multirow{3}{3.5cm}{Condiciones institucionales (factores asociados)} & 
1. Concentración discrecional de poder sin contrapesos. \newline
2. Debilidad en preinversión y fiscalización de obra pública. \newline
3. Desregulación y conflicto de interés en servicios subrogados. & 
Expansión operativa y presupuestal en la SSP sin procesamiento interno; auditorías técnicas de preinversión de la ASF; tasa de crecimiento de guarderías subrogadas frente a directas en el IMSS y vínculos familiares/políticos registrados en expedientes. \\
\addlinespace
\multirow{3}{3.5cm}{Efectos documentados (resultados observables)} & 
1. Captura institucional y macrocriminalidad. \newline
2. Daño patrimonial y pasivos fiscales transexenales. \newline
3. Vulneración de derechos humanos y deterioro de la confianza pública. & 
Sentencias penales y civiles en EE. UU. y México; montos desviados, porcentajes de sobrecosto y amortización de deuda estatal; víctimas mortales y lesionadas; trayectoria del IPC (2012--2025) y tasas de la ENCIG 2025. \\
\addlinespace
Variable de contraste (encuadre mediático) & 
Diferencia entre atención mediática (framing) y daño verificable a la población. & 
Contraste entre titulares virales no acreditados o reduccionistas y hallazgos oficiales de auditorías, verificaciones sanitarias (COFEPRIS) y reportes operativos (Veracruz y Dos Bocas). \\
\bottomrule
\end{tabular}
\end{table}